
\documentclass{article}
\usepackage{colortbl}
\usepackage{makecell}
\usepackage{multirow}
\usepackage{supertabular}

\begin{document}

\newcounter{utterance}

\centering \large Interaction Transcript for game `cladder', experiment `full\_v1.5\_default', episode 949 with reflexion{-}llama{-}3.1{-}8b{-}instruct.
\vspace{24pt}

{ \footnotesize  \setcounter{utterance}{1}
\setlength{\tabcolsep}{0pt}
\begin{supertabular}{c@{$\;$}|p{.15\linewidth}@{}p{.15\linewidth}p{.15\linewidth}p{.15\linewidth}p{.15\linewidth}p{.15\linewidth}}
    \# & \multicolumn{2}{c}{Player} && \multicolumn{2}{c}{Game Master} \\
    \hline

    \theutterance \stepcounter{utterance}  
    & & & \multicolumn{4}{p{0.6\linewidth}}{
        \cellcolor[rgb]{0.9,0.9,0.9}{
            \makecell[{{p{\linewidth}}}]{
                \texttt{\tiny{[P1$\langle$GM]}}
                \texttt{You are an expert in causal inference. The following question is not a typical commonsense query, but rather a meticulously designed question created by a professor specializing in causal inference, intended to assess the students' mastery of the course content.} \\
\texttt{Imagine a self{-}contained, hypothetical world with only the following conditions, and without any unmentioned factors or causal relationships: Obesity has a direct effect on diabetes and lifespan. Smoking has a direct effect on diabetes and lifespan. Diabetes has a direct effect on lifespan. Smoking is unobserved. The overall probability of obesity is 10\%. For normal weight people, the probability of long lifespan is 86\%. For obese people, the probability of long lifespan is 54\%. Is long lifespan less likely than short lifespan overall?} \\
\texttt{Start your answer with 'Yes' or 'No', followed by additional reasoning or evidence to support your explanation.} \\
            }
        }
    }
     \\ \\

    \theutterance \stepcounter{utterance}  
    & \multicolumn{4}{p{0.6\linewidth}}{
        \cellcolor[rgb]{0.561,0.737,0.561}{
            \makecell[{{p{\linewidth}}}]{
                \texttt{\tiny{[P1$\rangle$GM]}}
                \texttt{No, long lifespan is not less likely than short lifespan overall.} \\
\\ 
\texttt{To answer this question, we need to apply the principles of causal inference, specifically the concept of counterfactuals and potential outcomes. Since smoking is unobserved, we cannot directly estimate its effect on lifespan. However, we can use the information provided to make an indirect inference.} \\
\\ 
\texttt{Let's define the following variables:} \\
\\ 
\texttt{{-} O: Obesity (observed)} \\
\texttt{{-} D: Diabetes (observed)} \\
\texttt{{-} L: Long lifespan (observed)} \\
\texttt{{-} S: Smoking (unobserved)} \\
\\ 
\texttt{We are given the following probabilities:} \\
\\ 
\texttt{{-} P(O) = 0.1 (probability of obesity)} \\
\texttt{{-} P(L|O) = 0.54 (probability of long lifespan given obesity)} \\
\texttt{{-} P(L|not O) = 0.86 (probability of long lifespan given normal weight)} \\
\\ 
\texttt{We also know that diabetes has a direct effect on lifespan, and obesity and smoking have direct effects on diabetes and lifespan. However, since smoking is unobserved, we cannot directly estimate its effect on lifespan.} \\
\\ 
\texttt{To determine whether long lifespan is less likely than short lifespan overall, we need to consider the potential outcomes of being obese or normal weight. Let's define the following potential outcomes:} \\
\\ 
\texttt{{-} L(O): Long lifespan given obesity} \\
\texttt{{-} S(O): Short lifespan given obesity} \\
\texttt{{-} L(not O): Long lifespan given normal weight} \\
\texttt{{-} S(not O): Short lifespan given normal weight} \\
\\ 
\texttt{We can estimate the overall probability of long lifespan as follows:} \\
\\ 
\texttt{P(L) = P(L|O) \textbackslash{}* P(O) + P(L|not O) \textbackslash{}* P(not O)} \\
\texttt{= 0.54 \textbackslash{}* 0.1 + 0.86 \textbackslash{}* 0.9} \\
\texttt{= 0.774} \\
\\ 
\texttt{Similarly, we can estimate the overall probability of short lifespan as follows:} \\
\\ 
\texttt{P(S) = P(S|O) \textbackslash{}* P(O) + P(S|not O) \textbackslash{}* P(not O)} \\
\\ 
\texttt{However, we do not have direct information about P(S|O) or P(S|not O). To make an indirect inference, we can use the fact that diabetes has a direct effect on lifespan. Let's define the following variables:} \\
\\ 
\texttt{{-} D(O): Diabetes given obesity} \\
\texttt{{-} D(not O): Diabetes given normal weight} \\
\\ 
\texttt{We can estimate the probability of diabetes given obesity as follows:} \\
\\ 
\texttt{P(D|O) = P(D|O and S) \textbackslash{}* P(S|O) + P(D|O and not S) \textbackslash{}* P(not S|O)} \\
\texttt{= P(D|O and S) \textbackslash{}* P(S|O) + P(D|O and not S) \textbackslash{}* (1 {-} P(S|O))} \\
\\ 
\texttt{However, we do not have direct information about P(S|O) or P(D|O and S). To make an indirect inference, we can use the fact that obesity has a direct effect on diabetes. Let's define the following variables:} \\
\\ 
\texttt{{-} O(D):} \\
            }
        }
    }
    & & \\ \\

    \theutterance \stepcounter{utterance}  
    & & & \multicolumn{2}{p{0.3\linewidth}}{
        \cellcolor[rgb]{0.9,0.9,0.9}{
            \makecell[{{p{\linewidth}}}]{
                \texttt{\tiny{[GM$|$GM]}}
                \texttt{No} \\
            }
        }
    }
    & & \\ \\

    \theutterance \stepcounter{utterance}  
    & & & \multicolumn{2}{p{0.3\linewidth}}{
        \cellcolor[rgb]{0.9,0.9,0.9}{
            \makecell[{{p{\linewidth}}}]{
                \texttt{\tiny{[GM$|$GM]}}
                \texttt{no} \\
            }
        }
    }
    & & \\ \\

    \theutterance \stepcounter{utterance}  
    & & & \multicolumn{2}{p{0.3\linewidth}}{
        \cellcolor[rgb]{0.9,0.9,0.9}{
            \makecell[{{p{\linewidth}}}]{
                \texttt{\tiny{[GM$|$GM]}}
                \texttt{game\_result = WIN} \\
            }
        }
    }
    & & \\ \\

\end{supertabular}
}

\end{document}
